\subsection{Universal coefficient formula}

In this section, we consider two complexes of A-modules (C, d) and (C', d'). Consider the canonical exact sequences:

\[
0 \rightarrow \mathrm{Z} (\mathrm{C}) \stackrel {{j}} {{\rightarrow}} \mathrm{C} \stackrel {{\delta}} {{\rightarrow}} \mathrm{B} (\mathrm{C}) (- 1) \rightarrow 0,\tag{I}
\]

\((\mathrm{II}_{p})\)

\[
0 \to \mathrm{B} _ {p} (\mathrm{C}) \stackrel {{i}} {{\to}} \mathrm{Z} _ {p} (\mathrm{C}) \stackrel {{p}} {{\to}} \mathrm{H} _ {p} (\mathrm{C}) \to 0;
\]

we deduce from δ a k-homomorphism

\[
\mathrm{H} (\text { Homgr } (\delta , 1)): \mathrm{H} (\text { Homgr } _ {\mathbf {A}} (\mathrm{B} (\mathbf {C}), \mathbf {C} ^ {\prime})) (1) \rightarrow \mathrm{H} (\text { Homgr } _ {\mathbf {A}} (\mathbf {C}, \mathbf {C} ^ {\prime}));
\]

one deduces from \((\mathrm{II}_{p})\) connecting homomorphisms:

\[
\delta \left(\mathrm{II} _ {p}, \mathrm{H} ^ {q} \left(\mathrm{C} ^ {\prime}\right)\right): \operatorname{Hom} _ {\mathrm{A}} \left(\mathrm{B} _ {p} (\mathrm{C}), \mathrm{H} ^ {q} \left(\mathrm{C} ^ {\prime}\right)\right)\rightarrow \operatorname{Ext} _ {\mathrm{A}} ^ {1} \left(\mathrm{H} _ {p} (\mathrm{C}), \mathrm{H} ^ {q} \left(\mathrm{C} ^ {\prime}\right)\right).
\]

whence, by passing to the product, homomorphisms of k-modules

\[
\varphi^ {n}: \operatorname{Homgr} _ {\mathbf {A}} ^ {n} (\mathrm{B} (\mathrm{C}), \mathrm{H} (\mathrm{C} ^ {\prime})) \rightarrow \prod_ {p + q = n} \operatorname{Ext} _ {\mathbf {A}} ^ {1} \left(\mathrm{H} _ {p} (\mathrm{C}), \mathrm{H} ^ {q} (\mathrm{C} ^ {\prime})\right)
\]

One also has canonical homomorphisms (X, p. 82)

\[
\lambda^ {n} (\mathrm{B} (\mathrm{C}), \mathrm{C} ^ {\prime}): \mathrm{H} ^ {n} \left(\operatorname{Homgr} _ {\mathrm{A}} (\mathrm{B} (\mathrm{C}), \mathrm{C} ^ {\prime})\right)\rightarrow \operatorname{Homgr} _ {\mathrm{A}} ^ {n} (\mathrm{B} (\mathrm{C}), \mathrm{H} (\mathrm{C} ^ {\prime})) .
\]

With these notations:

THEOREM 3. --- Suppose the A-modules B(C) and Z(C) are projective. There exists, for each n, a unique homomorphism of k-modules

\[
\beta^ {n}: \prod_ {p + q = n - 1} \operatorname{Ext} _ {\mathbf {A}} ^ {1} \left(\mathrm{H} _ {p} (\mathbf {C}), \mathrm{H} ^ {q} \left(\mathbf {C} ^ {\prime}\right)\right)\rightarrow \mathrm{H} ^ {n} \left(\operatorname{Homgr} _ {\mathbf {A}} \left(\mathrm{C}, \mathrm{C} ^ {\prime}\right)\right)
\]

thereby making the diagram commutative

\[
\begin{array}{c} \operatorname{Homgr} _ {\mathbf {A}} ^ {n - 1} (\mathrm{B} (\mathrm{C}), \mathrm{H} (\mathrm{C} ^ {\prime})) \xleftarrow {\lambda^ {n - 1} (\mathrm{B} (\mathrm{C}) , \mathrm{C} ^ {\prime})} \mathrm{H} ^ {n - 1} (\operatorname{Homgr} _ {\mathbf {A}} (\mathrm{B} (\mathrm{C}), \mathrm{C} ^ {\prime})) \\ \prod_ {p + q = n - 1} \operatorname{Ext} _ {\mathbf {A}} ^ {1} (\mathrm{H} _ {p} (\mathrm{C}), \mathrm{H} ^ {q} (\dot {\mathrm{C}} ^ {\prime})) \xrightarrow {\beta^ {n}} \mathrm{H} ^ {n} (\operatorname{Homgr} _ {\mathbf {A}} (\mathrm{C}, \mathrm{C} ^ {\prime})). \end{array}
\]

The sequences of graded k-modules

\[
\begin{array}{r l} 0 \longrightarrow & \prod_ {p + q = n - 1} \operatorname{Ext} _ {\mathrm{A}} ^ {1} \left(\mathrm{H} _ {p} (\mathrm{C}), \mathrm{H} ^ {q} (\mathrm{C} ^ {\prime})\right) \xrightarrow {\beta^ {n}} \mathrm{H} ^ {n} \big (\mathrm{Homg} _ {\mathrm{A}} (\mathrm{C}, \mathrm{C} ^ {\prime}) \big) \\ & \xrightarrow {\lambda^ {n} (\mathrm{C} , \mathrm{C} ^ {\prime})} \prod_ {p + q = n} \mathrm{Hom} _ {\mathrm{A}} \left(\mathrm{H} _ {p} (\mathrm{C}), \mathrm{H} ^ {q} (\mathrm{C} ^ {\prime})\right) \to 0 \end{array}\tag{11}
\]

are exact.

Remark. --- One can prove an analogous statement, assuming \(\mathbf{B}_{n}(\mathbf{C}')\) and \(\mathbf{C}'/\mathbf{B}_{n}(\mathbf{C}')\) injective for each n.

For simplicity set \(\mathbf{B} = \mathbf{B}(\mathbf{C}), \mathbf{Z} = \mathbf{Z}(\mathbf{C}), \mathbf{H} = \mathbf{H}(\mathbf{C})\) and \(\mathbf{H}' = \mathbf{H}(\mathbf{C}')\) . Since B is projective, from (I) one derives an exact sequence

\[
\begin{array}{r l}(1 2) \quad&0 \rightarrow \operatorname{Homgr} _ {\mathbf {A}} (\mathrm{B}, \mathrm{C} ^ {\prime}) (1) \xrightarrow {\operatorname{Homgr} (\delta , 1)} \operatorname{Homgr} _ {\mathbf {A}} (\mathrm{C}, \mathrm{C} ^ {\prime})\\&\xrightarrow {\operatorname{Homgr} (j , 1)} \operatorname{Homgr} _ {\mathbf {A}} (\mathrm{Z}, \mathrm{C} ^ {\prime}) \rightarrow 0.\end{array}
\]

Lemma 2. --- The connecting homomorphism

\[
\mathrm{H} ^ {n} \left(\operatorname{Homgr} _ {\mathbf {A}} \left(\mathrm{Z}, \mathrm{C} ^ {\prime}\right)\right)\rightarrow \mathrm{H} ^ {n} \left(\operatorname{Homgr} _ {\mathbf {A}} \left(\mathrm{B}, \mathrm{C} ^ {\prime}\right)\right)
\]

associated with (12) is equal to \((-1)^{n+1}\) H(Homgr(i, 1)).

Indeed, let \(a \in \mathbf{Z}^n(\text{Homgr}_A(Z, C'))\) ; this is a morphism of complexes of ascending degree \(n\) of \(Z\) to \(C'\) , whose values therefore lie in \(Z(C')\) . Since the exact sequence (I) is split (B being projective), \(a\) extends to an element \(b\) of \(\text{Homgr}_A^n(C, Z(C'))\) . By definition, the image of the class of \(a\) under the connecting homomorphism in question is the class in \(H^n(\text{Homgr}_A(B, C'))\) of the homomorphism \(u\) of \(B(C)\) to \(C'\) such that for \(x \in C\) , one has

\[
u (d x) = \mathrm{D} b (x) = d ^ {\prime} b (x) - (- 1) ^ {n} b (d x) = (- 1) ^ {n + 1} b (d x) = (- 1) ^ {n + 1} a (d x),
\]

whence the assertion.

The exact homology sequence associated with (12) therefore gives the exact sequence

\[
\rightarrow \mathrm{H} ^ {n} (\operatorname{Homgr} _ {\mathrm{A}} (Z, C ^ {\prime})) \xrightarrow {\mathrm{H} (\operatorname{Homgr} (i , 1))} \mathrm{H} ^ {n} (\operatorname{Homgr} _ {\mathrm{A}} (B, C ^ {\prime}))
\]

\[
\xrightarrow {\mathrm{H} (\text { Homgr } (\delta , 1))} \mathrm{H} ^ {n + 1} \left(\text { Homgr } _ {\mathrm{A}} \left(\mathrm{C}, \mathrm{C} ^ {\prime}\right)\right) \xrightarrow {\mathrm{H} (\text { Homgr } (j , 1))} \mathrm{H} ^ {n + 1} \left(\text { Homgr } _ {\mathrm{A}} \left(\mathrm{Z}, \mathrm{C} ^ {\prime}\right)\right)\rightarrow \dots .
\]

Moreover, since Z is projective, one derives from \((II_{p})\) exact sequences

\[
0 \to \mathrm{Hom} _ {\mathbf {A}} (\mathrm{H} _ {p}, \mathrm{H} ^ {\prime q}) \to \mathrm{Hom} _ {\mathbf {A}} (Z _ {p}, \mathrm{H} ^ {\prime q}) \to \mathrm{Hom} _ {\mathbf {A}} (\mathrm{B} _ {p}, \mathrm{H} ^ {\prime q}) \to \mathrm{Ext} _ {\mathbf {A}} ^ {1} (\mathrm{H} _ {p}, \mathrm{H} ^ {\prime q}) \to 0,
\]

whence, by passing to products, exact sequences

\[
\begin{array}{r l} 0 \to \operatorname{Homgr} _ {\mathbf {A}} ^ {n} (\mathbf {H}, \mathbf {H} ^ {\prime}) & \to \operatorname{Homgr} _ {\mathbf {A}} ^ {n} (Z, \mathbf {H} ^ {\prime}) \to \operatorname{Homgr} _ {\mathbf {A}} ^ {n} (\mathbf {B}, \mathbf {H} ^ {\prime}) \\ & \to \prod_ {p + q = n} \operatorname{Ext} _ {\mathbf {A}} ^ {1} (\mathbf {H} _ {p}, \mathbf {H} ^ {\prime q}) \to 0. \end{array}
\]

Finally, one has the canonical homomorphisms of no. 1:

\[
\begin{array}{l} \lambda_ {\mathrm{B}} = \lambda (\mathrm{B}, \mathrm{C} ^ {\prime}): \mathrm{H} (\mathrm{Homgr} _ {\mathrm{A}} (\mathrm{B}, \mathrm{C} ^ {\prime})) \to \mathrm{Homgr} _ {\mathrm{A}} (\mathrm{B}, \mathrm{H} ^ {\prime}), \\ \lambda_ {\mathrm{Z}} = \lambda (\mathrm{Z}, \mathrm{C} ^ {\prime}): \mathrm{H} (\mathrm{Homgr} _ {\mathrm{A}} (\mathrm{Z}, \mathrm{C} ^ {\prime})) \to \mathrm{Homgr} _ {\mathrm{A}} (\mathrm{Z}, \mathrm{H} ^ {\prime}), \\ \lambda_ {\mathrm{C}} = \lambda (\mathrm{C}, \mathrm{C} ^ {\prime}): \mathrm{H} (\mathrm{Homgr} _ {\mathrm{A}} (\mathrm{C}, \mathrm{C} ^ {\prime})) \to \mathrm{Homgr} _ {\mathrm{A}} (\mathrm{H}, \mathrm{H} ^ {\prime}), \end{array}
\]

whence the diagram with exact rows on page X. 97.

This diagram is commutative by construction of the homomorphisms \(\lambda\) . Moreover, since the complexes B and Z are split and projective, \(\lambda_{\mathrm{B}}\) and \(\lambda_{\mathrm{Z}}\) are bijective (X, p. 84, cor. 1). From this one deduces, on the one hand, that \(\lambda_{\mathrm{C}}^{n}\) is surjective, with kernel equal to Im H \(^{n}\) (Homgr ( \(\delta\) , 1)), and on the other hand, that \(\varphi^{n-1} \circ \lambda_{\mathrm{B}}^{n-1}\) is surjective, with kernel equal to Ker H \(^{n}\) (Homgr ( \(\delta\) , 1)). The theorem follows immediately from this.

COROLLARY 1. --- Suppose B(C) and Z(C) are projective and B \(^{n}\) (C') is injective for each n. Then the exact sequences (11) are split.

This follows from the theorem and the following lemma:

Lemma 3. --- If B(C) is projective and \(\mathbf{B}_{n}(\mathbf{C}^{\prime})\) injective for each \(n\) , the canonical homomorphism \(\lambda(\mathbf{C}, \mathbf{C}^{\prime}): \mathrm{H}(\mathrm{Homgr}_{\mathbf{A}}(\mathbf{C}, \mathbf{C}^{\prime})) \to \mathrm{Homgr}_{\mathbf{A}}(\mathrm{H}(\mathbf{C}), \mathrm{H}(\mathbf{C}^{\prime}))\) has a section \(k\) -linear.

Indeed, according to X, p. 35, remarks a) and b), there exist homologisms

\[
\varphi : \mathrm{C} \rightarrow \mathrm{H} (\mathrm{C}) \quad \text { and } \quad \varphi^ {\prime}: \mathrm{H} (\mathrm{C} ^ {\prime}) \rightarrow \mathrm{C} ^ {\prime}
\]

such that \(\mathbf{H}(\varphi) = 1_{\mathbf{H}(\mathbf{C})}\) and \(\bar{\mathbf{H}} (\varphi^{\prime}) = 1_{\mathbf{H}(\mathbf{C}^{\prime})}\) .

In the commutative diagram

\[
\begin{array}{l} \mathrm{H} (\operatorname{Homgr} _ {\mathbf {A}} (\mathrm{H} (\mathrm{C}), \mathrm{H} (\mathrm{C} ^ {\prime})) \xrightarrow {\mathrm{H} (\operatorname{Homgr} (\varphi , \varphi^ {\prime}))} \mathrm{H} (\operatorname{Homgr} _ {\mathbf {A}} (\mathrm{C}, \mathrm{C} ^ {\prime})) \\ \lambda (\mathrm{H} (\mathrm{C}), \mathrm{H} (\mathrm{C} ^ {\prime})) \Big | _ {\downarrow} \\ \operatorname{Homgr} _ {\mathbf {A}} (\mathrm{H} (\mathrm{C}), \mathrm{H} (\mathrm{C} ^ {\prime})) \xrightarrow {\operatorname{Homgr} (\mathrm{H} (\varphi) , \mathrm{H} (\varphi^ {\prime}))} \operatorname{Homgr} _ {\mathbf {A}} (\mathrm{H} (\mathrm{C}), \mathrm{H} (\mathrm{C} ^ {\prime}))  , \end{array}
\]

\(\lambda(H(C), H(C'))\) is bijective and Homgr \((H(\varphi), H(\varphi'))\) is the identity, whence the assertion.

Homgr \(^{n-1}\) (Z, H') → Homgr \(^{n-1}\) (B, H') → Φ \(^{n-1}\) → Πp+q=n-1 Ext \(^{1}\) (Hp, H'q) → 0

H \(^{n-1}\) (Homgr(A, C')) → H \(^{n-1}\) (Homgr(A, B, C')) → H \(^{n}\) (Homgr(A, C)) → H \(^{n}\) (Homgr(A, C)) → H \(^{n}\) (Homgr(A, B, C))

0 → Homgr \(^{n}\) (H, H') → Homgr \(^{n}\) (Z, H') → Homgr \(^{n}\) (B, H')

COROLLARY 2. --- If B(C) and Z(C) are projective, then, for every A-module N and every integer n, there is a split exact sequence:

\[
(1 3) \quad 0 \longrightarrow \operatorname{Ext} _ {\mathbf {A}} ^ {1} \left(\mathrm{H} _ {n - 1} (\mathrm{C}), \mathrm{N}\right) \xrightarrow {\beta^ {n}} \mathrm{H} ^ {n} \left(\operatorname{Homgr} _ {\mathbf {A}} (\mathrm{C}, \mathrm{N})\right) \xrightarrow {\lambda^ {n}} \operatorname{Hom} _ {\mathbf {A}} \left(\mathrm{H} _ {n} (\mathrm{C}), \mathrm{N}\right) \longrightarrow 0.
\]

COROLLARY 3 (“universal coefficient formula”). --- Suppose A is principal and C is free. Then for every A-module N and every integer n, there is a split exact sequence (13).

Indeed, B(C) and Z(C) are free as submodules of the free module C (VII, § 3, cor. 2 to th. 1).

COROLLARY 4. --- If C is bounded on the right and if C and H(C) are projective, then

\[
\lambda (C, C ^ {\prime}): H \left(\operatorname{Homgr} _ {A} (C, C ^ {\prime})\right)\rightarrow \operatorname{Homgr} _ {A} \left(H (C), H \left(C ^ {\prime}\right)\right)
\]

is bijective.

By the theorem, it suffices to prove that B(C) and Z(C) are projective. Now we have exact sequences

\[
\begin{array}{l} 0 \to \mathrm{B} _ {n} (\mathrm{C}) \to Z _ {n} (\mathrm{C}) \to \mathrm{H} _ {n} (\mathrm{C}) \to 0 \\ 0 \to Z _ {n} (\mathrm{C}) \to \mathrm{C} _ {n} \to \mathrm{B} _ {n - 1} (\mathrm{C}) \to 0, \end{array}
\]

hence \((\mathbf{B}_{n - 1}(\mathbf{C})\) is projective) \(\Rightarrow (\mathbf{Z}_n(\mathbf{C})\) is projective) \(\Rightarrow (\mathbf{B}_n(\mathbf{C})\) is projective). We conclude by noting that \(\mathbf{B}_n(\mathbf{C}) = 0\) for \(n\) small enough.

